\documentclass[aps,prb,twocolumn,superscriptaddress,floatfix,nofootinbib]{revtex4-2}

% ----------- Paquetes básicos -----------
\usepackage[utf8]{inputenc}
\usepackage[T1]{fontenc}
\usepackage{amsmath,amssymb,mathtools}
\usepackage{siunitx}
\usepackage{graphicx}
\usepackage{booktabs}
\usepackage[dvipsnames]{xcolor}
\usepackage{hyperref}
\hypersetup{
  colorlinks=true,
  linkcolor=MidnightBlue,
  citecolor=MidnightBlue,
  urlcolor=MidnightBlue
}

% ----------- Notación -----------
\newcommand{\Tc}{\ensuremath{T_{\mathrm{c}}}}
\newcommand{\Tcone}{\ensuremath{{T_{\mathrm{c}}^{\mathrm{ideal}}}}} 
\newcommand{\ThetaD}{\ensuremath{\Theta_{\mathrm{D}}}}
\newcommand{\Fm}{\ensuremath{F_{\mathrm{m}}}}
\newcommand{\Fmstar}{\ensuremath{{F_{\mathrm{m}}^{\ast}}}} 
\newcommand{\fbase}{\ensuremath{{f_{\mathrm{base}}}}} 
\newcommand{\Ncorr}{\ensuremath{{N_{\mathrm{corr}}}}}
\newcommand{\deltalock}{\ensuremath{\delta}}
\newcommand{\znoise}{\ensuremath{{z_{\xi}}}}

\newcommand{\He}{\ensuremath{\mathrm{He}}}
\newcommand{\HeFour}{\ensuremath{{^{4}\He}}}

\newcommand{\figref}[1]{Fig.~\ref{#1}}
\newcommand{\secref}[1]{Sec.~\ref{#1}}
% ----------- Documento -----------
\begin{document}

\title{Integer participation and structural noise in superconducting clusters}

\author{C.~Agostino}
\affiliation{Independent Researcher}

\date{\today}

% ----------- Abstract -----------
\begin{abstract}
We analyze a curated dataset of more than two hundred superconducting materials
and liquid \HeFour\ using a single phenomenological observable,
the ``mother frequency'' $\Fm = \ThetaD / \Tc$,
where $\ThetaD$ is an effective Debye temperature and $\Tc$ the transition temperature.
Using a structural--noise model inferred from previous work on cluster topology, frictional heating and pressure sensitivity, we define an
ideal transition temperature
$\Tcone = \Tc (1 + \widehat{\xi})$ and a noise--corrected
frequency $\Fmstar = \ThetaD / \Tcone$.
We then search, for each family of materials, a base frequency
$\fbase$ that minimizes the median distance of
$N_{\mathrm{corr}} = \Fmstar / \fbase$ to the nearest integer.
The resulting integer ``participation numbers'' $N_{\mathrm{corr}}$ exhibit
a strong deviation from several null models
obtained by shuffling $\ThetaD$ or sampling from smooth continuous distributions.
The distribution of $|\deltalock| = |\Ncorr - \mathrm{round}(\Ncorr)|$
is significantly more concentrated near zero than the nulls
(Kolmogorov--Smirnov and Mann--Whitney $p \lesssim 10^{-13}$,
Cliff's $\Delta \simeq 0.3$),
indicating genuine integer locking rather than a discretization artefact.
Across families, we find a robust hierarchy of base frequencies:
iron-based superconductors cluster near $\fbase \simeq 1$,
effectively behaving as natural, unscaled harmonic oscillators for which
$\Fmstar$ itself is already an (almost) integer-valued participation number;
binary, molecular and high-pressure materials sit near $\fbase \simeq 4$,
and oxides and heavy fermions around $\fbase \simeq 2$.
The global median $\fbase \simeq 4.16$ is numerically close
to the $n=4$ thermal ratio previously identified between the
superfluid and roton branches of liquid \HeFour.
Finally, we show that materials with almost-integer participation
systematically exhibit lower structural noise:
the top 20\% with smallest $|\deltalock|$ have a substantially
lower noise $z$--score than the rest of the dataset
(Cliff's $\Delta \simeq -0.47$).
Together, these results suggest that the mother frequency $\Fm$
probes a discrete oscillator structure underpinning both superconductivity
and superfluidity, with family-dependent harmonics shaped
by structural disorder and pressure.
\end{abstract}

\maketitle

% ----------- I. INTRODUCTION -----------
\section{Introduction}
\label{sec:intro}

The search for organizing principles in the zoo of superconducting
materials typically proceeds through microscopic models for pairing,
band structure or electron--phonon coupling.\cite{BCS1957,Tinkham1996}
Here we explore a complementary, purely phenomenological direction,
inspired by methods in non-equilibrium statistical mechanics and
dynamical systems~\cite{Mori1965,Zwanzig2001,Pikovsky2001,GuckenheimerHolmes1983}:
we treat each superconductor or superfluid as a black box characterized
by a single dimensionless frequency
\begin{equation}
  \Fm = \frac{\ThetaD}{\Tc},
\end{equation}
constructed from an effective Debye temperature $\ThetaD$ and
the observed critical temperature $\Tc$.
Previous work showed that, for liquid \HeFour, ratios of such
frequencies across well-defined dynamical regimes (normal fluid, superfluid, roton)
fall extremely close to small integers, with a particularly prominent
factor $n = 4$ associated with the onset of the roton branch.
This suggests that, at least for \HeFour, the spectrum of collective excitations
organizes into discrete harmonic ``shells''.

In this paper we ask whether an analogous discrete structure can be detected,
in a data-driven way, across a heterogeneous ensemble of superconductors.
Rather than committing to a specific microscopic Hamiltonian,
we use $\Fm$ as a coarse-grained probe and search for hidden
integer structure once known sources of noise are removed.

Two ingredients are essential:
(i) a structural--noise model that accounts for deviations of the
measured $\Tc$ from an ideal value set by frictional losses, topology
and pressure; and
(ii) a robust calibration of a base frequency $\fbase$ that captures
the dominant periodicity of $\Fm$ within each family.
After correcting for noise, we show that the rescaled quantities
$\Ncorr = \Fmstar / \fbase$ cluster close to integers
far more strongly than expected from smooth null models.
We further identify a family-dependent hierarchy of
$\fbase$ that lines up convincingly with the integer ratios
previously observed in \HeFour.

The rest of this paper is organized as follows.
\secref{sec:data} describes the dataset and the structural--noise model.
\secref{sec:methods} defines the base-frequency calibration
and the integer-participation metrics.
\secref{sec:results} presents the statistical evidence for
integer locking and its dependence on structural noise and material family.
\secref{sec:discussion} discusses physical interpretations and limitations,
and \secref{sec:conclusion} summarizes the main conclusions.
Additional figures validating the noise model are provided in the Appendix.

% ----------- II. DATA AND NOISE MODEL -----------
\section{Data and structural--noise model}
\label{sec:data}

\subsection{Dataset and family labels}

We work with the same curated dataset of superconducting materials
used in previous studies of cluster topology and pressure sensitivity.
Each entry corresponds to a material or fluid labeled by a name and
a category: \texttt{SC\_TypeI}, \texttt{SC\_TypeII}, \texttt{SC\_Binary},
\texttt{SC\_Oxide}, \texttt{SC\_Molecular}, \texttt{SC\_HighPressure},
\texttt{SC\_HeavyFermion}, \texttt{SC\_IronBased}, and \texttt{Superfluid}
(for liquid \HeFour).
For each entry $i$ we have an effective Debye temperature $\ThetaD^{(i)}$,
a measured transition temperature $\Tc^{(i)}$,
a family label $k(i)$, and a structural--noise prediction
$\widehat{\xi}^{(i)}$ obtained from a separate model
that encodes frictional dissipation, topological defects
and pressure-induced distortions.

\subsection{Structural noise and ideal transition temperature}

The structural--noise field $\widehat{\xi}^{(i)}$ is defined such that
$\widehat{\xi}^{(i)} \ge 0$ corresponds to a loss of effective ordering
relative to an otherwise similar but better structured material.
We therefore interpret it as a multiplicative distortion of the
transition temperature and define an ``ideal'' transition
\begin{equation}
  \label{eq:Tc_ideal}
  \Tcone^{(i)} = \Tc^{(i)} \bigl[1 + \widehat{\xi}^{(i)}\bigr].
\end{equation}
By construction, $\Tcone^{(i)} \ge \Tc^{(i)}$ and the correction vanishes
for structurally optimal samples.

The corresponding noise-corrected mother frequency is
\begin{equation}
  \label{eq:Fm_star}
  \Fmstar^{(i)} = \frac{\ThetaD^{(i)}}{\Tcone^{(i)}}.
\end{equation}
\figref{fig:Tc_vs_Tc_ideal} (Appendix) confirms that the relation
between $\Tc$ and $\Tcone$ is close to the identity for most materials,
with a small subset receiving sizable upward corrections.

To compare noise levels across families with different baselines,
we standardize the predicted noise within each family $k$:
\begin{equation}
  \label{eq:z_noise}
  \znoise^{(i)} =
  \frac{\widehat{\xi}^{(i)} - \mu_{\xi,k(i)}}{\sigma_{\xi,k(i)}},
\end{equation}
where $\mu_{\xi,k}$ and $\sigma_{\xi,k}$ are the mean and
standard deviation of $\widehat{\xi}$ in family $k$.
The resulting distributions of $\znoise$ are shown in
\figref{fig:noise_by_family} (Appendix).

% -- TO APPROVE --
\subsection{Validation against null models}\label{sec:integer_validation}

Because \(f_{\mathrm{base},k}\) is determined by minimising an integer loss,
it is essential to test whether the resulting structure can be explained by
overfitting or by generic features of positive random variables. We therefore
compare the empirical loss \(\tilde L_k\) with several null models.

First, for each family we construct shuffled datasets by randomly permuting
the Debye temperatures \(\Theta_D\) across materials while keeping
\(T_{c,\mathrm{ideal}}\) fixed. For each permutation we recompute
\(F_m^*\), refit \(f_{\mathrm{base},k}\), and obtain a shuffled loss
\(\tilde L_k^{\mathrm{shuff}}\). The fraction of permutations with
\(\tilde L_k^{\mathrm{shuff}} \le \tilde L_k\) defines a non-parametric
\(p\)-value for the integer structure.

Second, we generate synthetic Mother Frequencies from continuous
distributions (gamma or log-normal) fitted to the empirical \(F_m^*\) in
each family, again refitting \(f_{\mathrm{base},k}\) for each synthetic
realisation. This tests whether the observed integer alignment can be
reproduced purely from the one-point statistics of \(F_m^*\).

Finally, we quantify the link between integer participation and structural
noise by computing the Spearman rank correlation between the integer
residual
\(\delta_i = |N_i - \mathrm{round}(N_i)|\) and the normalised noise level
\(Z_{\xi,i}\) within and across families. We also compare the noise
distributions of ``almost-integer'' materials (e.g.\ the lowest 20\% in
\(\delta_i\)) against the rest using Cliff’s delta. In all cases we report
confidence intervals obtained by non-parametric bootstrap.

These tests show that the observed combination of small \(\tilde L_k\), low
integer residuals, and negative correlation between \(\delta_i\) and
\(Z_{\xi,i}\) is highly unlikely to arise from smooth random fields. The
integer structure therefore encodes information that is not already captured
by the scalar distribution of \(F_m^*\) or by the structural-noise field
alone.

\begin{figure}[t]
  \centering
  \includegraphics[width=\columnwidth]{figures/fig_delta_vs_zxi.png}
  \caption{Integer mismatch $\delta$ versus normalised structural noise $Z_\xi$
  across all families.
  Each point corresponds to a material; colours denote the superconducting
  family.
  The positive Spearman rank correlation
  ($\rho \approx 0.34$, $p \approx 5.6\times10^{-4}$) indicates that materials
  further from an integer participation number tend to have higher predicted
  structural noise, consistent with the structural-noise hypothesis.}
  \label{fig:delta_vs_zxi}
\end{figure}

% ----------- III. METHODS: BASE FREQ & INTEGER PARTICIPATION -----------
\section{Base-frequency calibration and integer participation}
\label{sec:methods}

\subsection{Base frequency per family}

We seek, for each family $k$, a base frequency $\fbase^{(k)}$
such that the rescaled frequencies
\begin{equation}
  \Ncorr^{(i)} = \frac{\Fmstar^{(i)}}{\fbase^{(k(i))}}
\end{equation}
cluster near integers.
To quantify this we define the integer locking residual
\begin{equation}
  \deltalock^{(i)} =
  \Ncorr^{(i)} - \mathrm{round}\!\bigl(\Ncorr^{(i)}\bigr),
\end{equation}
and focus on its absolute value $|\deltalock^{(i)}|$.

For each family $k$ we choose $\fbase^{(k)}$ by minimizing
a robust loss
\begin{equation}
  \label{eq:loss}
  \begin{split}
    \mathcal{L}_k(\fbase) = & \;\mathrm{median}_{i \in k}
    \Bigl|\Ncorr^{(i)} - \mathrm{round}\!\bigl(\Ncorr^{(i)}\bigr)\Bigr| \\
    & + \lambda \,\mathrm{median}_{i \in k} \Ncorr^{(i)},
  \end{split}
\end{equation}
over a physically motivated range
$0.5 \le \fbase \le 5.0$.
This explicitly rules out unphysical, high-frequency solutions
(e.g., $\fbase \sim 10^2$--$10^3$) that would artificially fit
higher harmonics while providing no interpretable physical scale.
The first term in Eq.~\eqref{eq:loss} drives the residuals toward integers;
the second is a small regularizer that penalizes large values of
$\Ncorr$ and selects low-order harmonics when multiple choices of
$\fbase$ give similar residuals.
The optimization is performed on a dense grid over the allowed range.
We repeat the calibration under two hypotheses,
using either $\Fmstar$ or $\Fmstar/2$ as the target frequency,
and keep the one with lower median residual.
All results reported below use this per-family calibration.

\begin{figure}[t]
  \centering
  \includegraphics[width=\columnwidth]{figures/fig_Lk_vs_f.png}
  \caption{Penalised loss $\tilde L(f)$ as a function of base frequency $f$
  for two representative families (Binary and High-pressure).
  Each curve is obtained by scanning $f$ in the range $[0.5,5.0]$ and computing
  the median integer mismatch plus a weak penalty proportional to the mean
  participation number.
  Dashed vertical lines mark the selected $f_{\mathrm{base}}$ values.
  The broad minima indicate that the calibration is not fine-tuned to a single
  spike in $f$.}
  \label{fig:loss_by_family}
\end{figure}

\begin{table*}[t]
  \caption{Summary of the superconducting dataset used in this work.
  $N_{\mathrm{mat}}$ is the number of distinct materials per family,
  $N_{\mathrm{rows}}$ the number of rows in the working table
  (including split sub-networks),
  $N_{\mathrm{single}}$, $N_{\sigma}$, $N_{\pi}$ and
  $N_{\mathrm{press}}$ count single, $\sigma$, $\pi$ and
  pressure--split networks, respectively.
  $T_c^{\mathrm{med}}$ and $\Theta_D^{\mathrm{med}}$ are the median
  critical and Debye temperatures, $\langle N\rangle$ is the mean
  participation number, and $\tilde\delta_{\mathrm{med}}$ the median
  integer mismatch.}
  \label{tab:dataset_summary}
  \begin{ruledtabular}
  \begin{tabular}{lrrrrrrrrrr}
    Family &
    $N_{\mathrm{mat}}$ &
    $N_{\mathrm{rows}}$ &
    $N_{\mathrm{single}}$ &
    $N_{\sigma}$ &
    $N_{\pi}$ &
    $N_{\mathrm{press}}$ &
    $T_c^{\mathrm{med}}$ (K) &
    $\Theta_D^{\mathrm{med}}$ (K) &
    $\langle N\rangle$ &
    $\tilde\delta_{\mathrm{med}}$ \\
    \hline
    SC\_Binary       & 67 & 77 & 57 & 3 & 3 & 0 &  6.00  & 280.0 & 155.4         & 0.175 \\
    SC\_HeavyFermion & 12 & 12 & 12 & 0 & 0 & 0 &  1.65  & 275.0 &  98.9         & 0.071 \\
    SC\_HighPressure & 80 & 82 & 79 & 2 & 1 & 0 & 145.5  & 925.0 &   1.95        & 0.180 \\
    SC\_IronBased    & 20 & 36 &  8 & 1 & 1 & 0 & 25.65  & 250.0 &  11.5         & 0.068 \\
    SC\_Molecular    & 10 & 10 & 10 & 0 & 0 & 0 & 11.00  & 145.0 &  13.5         & 0.090 \\
    SC\_Oxide        & 15 & 18 & 13 & 0 & 0 & 0 & 27.00  & 352.5 &  47.2         & 0.069 \\
    SC\_TypeI        & 10 & 11 &  9 & 0 & 0 & 0 &  3.20  & 200.0 &  28.2         & 0.153 \\
    SC\_TypeII       & 18 & 18 & 18 & 0 & 0 & 0 &  0.79  & 327.0 & $3.33\times10^{4}$ & 0.156 \\
    Superfluid       &  2 &  4 &  0 & 0 & 0 & 4 &  1.076 &  6.57 & 216.1         & 0.143 \\
  \end{tabular}
  \end{ruledtabular}
\end{table*}

\subsection{Null models}

To assess statistical significance, we compare the observed
locking residuals $|\deltalock^{(i)}|$ with several null models
that destroy any discrete structure while preserving
broad distributional features of $\Fmstar$~\cite{Efron1993}.

The simplest null shuffles the pairing between $\ThetaD$ and $\Tc$:
we randomly permute $\ThetaD$ across the dataset, recompute
$\Fm$ and $\Fmstar$ using Eqs.~\eqref{eq:Tc_ideal}--\eqref{eq:Fm_star},
recalibrate $\fbase^{(k)}$ and recompute $\Ncorr$ and $|\deltalock|$.
This preserves the marginal distributions of $\ThetaD$ and $\Tc$
and the family labels, but destroys any material-wise correlation.

We also consider smooth parametric nulls in which $\Fmstar$
is sampled from a fitted gamma or log-normal distribution within
each family, again followed by the same calibration procedure.
In all cases we generate $10^3$ realizations and pool the resulting
$|\deltalock|$ values.

% -- TO APPROVE --
\subsection{Structural noise field}\label{sec:structural_noise}

The integer–participation analysis in Sec.~\ref{sec:integer_participation}
requires a notion of ``structural noise'' for each material, encoding how far
its superconducting cluster departs from an ideal, stress–free configuration.
Rather than introducing this as a free parameter, we take \(\hat\xi\) from an
independent structural model.

Operationally, each material is assigned to a cluster topology described by
three discrete attributes: (i) macroscopic family (binary, type-I, type-II,
iron-based, high-pressure, etc.), (ii) sub-network (\textit{single}, \(\sigma\),
\(\pi\), or specialised high-pressure modes), and (iii) an effective
dimensionality \(d\in\{1,2,3\}\) that distinguishes quasi-one-dimensional
chains, layered systems, and three-dimensional networks. For this discrete
state space we construct a coarse structural noise field \(\hat\xi\) by
fitting a minimal regression model to a set of synthetic experiments where
cluster geometry, pressure sensitivity, and vector dissipation are explored
systematically (Study~02). The model and its calibration data are released as
CSV files and Python code in the accompanying repository, so that
\(\hat\xi\) can be recomputed or modified without re-fitting the DOFT
parameters.

In the present work we only require the sign and relative magnitude of the
structural noise. We therefore define an ``ideal'' critical temperature
\begin{equation}
  T_{c,\mathrm{ideal}} = T_c \bigl(1 + \hat\xi\bigr),
  \label{eq:Tc_ideal_def}
\end{equation}
which can be interpreted as the temperature that the material would achieve
if its cluster were fully relaxed with respect to the structural mode
captured by \(\hat\xi\). The corresponding renormalised Mother Frequency is
\begin{equation}
  F_m^* = \frac{\Theta_D}{T_{c,\mathrm{ideal}}}.
  \label{eq:Mother_freq_renorm}
\end{equation}

\begin{figure}[t]
  \centering
  \includegraphics[width=\columnwidth]{figures/figS02_Tc_vs_Tc_ideal.png}
  \caption{Comparison between reported critical temperatures $T_c$ and
  structurally corrected temperatures
  $T_c^{\mathrm{ideal}} = T_c(1+\hat\xi)$.
  The dashed line indicates $T_c^{\mathrm{ideal}} = T_c$; most points lie close
  to this line, with a few high-pressure hydrides shifted upwards by the
  structural-noise correction.}
  \label{fig:Tc_vs_Tcideal}
\end{figure}

Throughout the integer–participation analysis we work with \(F_m^*\) rather
than the raw ratio \(\Theta_D/T_c\); in Sec.~\ref{sec:integer_validation} we
show that this renormalisation improves the integer structure in a way that
cannot be reproduced by purely random perturbations.

\begin{figure*}[t]
  \centering
  \includegraphics[width=\textwidth]{figures/figS01_noise_by_family.png}
  \caption{Normalised structural noise $Z_\xi$ by family.
  Each box summarises the distribution of $Z_\xi$ for one superconducting family,
  after centering and scaling the predicted structural noise $\hat\xi$ within
  that family. Families are ordered by the median $Z_\xi$, so that groups on the
  left tend to exhibit lower predicted structural noise.}
  \label{fig:noise_by_family}
\end{figure*}

\subsection{Effect-size and correlation metrics}

We quantify deviations between the real and null distributions of
$|\deltalock|$ using non-parametric tests:
the Kolmogorov--Smirnov (KS) statistic,
the Mann--Whitney (MW) rank-sum test,
and Cliff's $\Delta$ as a measure of effect size~\cite{Cliff1993}.
Cliff's $\Delta$ is defined as
\begin{equation}
  \Delta =
  \Pr\bigl(|\deltalock|_{\mathrm{real}} <
           |\deltalock|_{\mathrm{null}}\bigr)
  - \Pr\bigl(|\deltalock|_{\mathrm{real}} >
             |\deltalock|_{\mathrm{null}}\bigr),
\end{equation}
and ranges from $-1$ to $+1$.

To connect locking strength with structural noise we compute,
over all materials with $\Ncorr^{(i)} \ge 1$,
the Spearman rank correlation between $|\deltalock^{(i)}|$
and $\znoise^{(i)}$.
We also split the dataset into the ``almost-integer'' group
(top 20\% lowest $|\deltalock|$) and the rest,
and compare their $\znoise$ distributions using Cliff's $\Delta$.

% -- TO APPROVE --
\subsection{Integer participation and base frequency}\label{sec:integer_participation}

The renormalised Mother Frequency \(F_m^*\) defined in
Eq.~\eqref{eq:Mother_freq_renorm} is dimensionless, but can be reinterpreted
as an effective participation number once a base frequency \(f_{\mathrm{base}}\)
is specified:
\begin{equation}
  N_i = \frac{F_{m,i}^*}{f_{\mathrm{base}}},
  \label{eq:integer_participation}
\end{equation}
where the index \(i\) runs over materials (or material–sub-network pairs).
The working hypothesis is that clean superconductivity corresponds to
near-integer values of \(N_i\), reflecting the coherent participation of a
discrete number of oscillatory units in the cluster.

We treat \(f_{\mathrm{base}}\) as a family-level parameter: for each
macroscopic family \(k\) we determine a separate \(f_{\mathrm{base},k}\) by
minimising the median distance to the nearest integer,
\begin{equation}
  L_k(f) = \mathrm{median}_i
           \bigl| N_i(f) - \mathrm{round}\bigl[N_i(f)\bigr] \bigr|,
  \qquad
  N_i(f) = \frac{F_{m,i}^*}{f},
  \label{eq:L_def}
\end{equation}
over a conservative interval \(f\in[f_{\min},f_{\max}]\). Using the median
rather than the mean makes the calibration robust to a small number of
anomalous materials, and the relatively narrow search range limits the
effective degrees of freedom. To avoid selecting artificially small base
frequencies that mimic higher harmonics, we add a weak penalty proportional
to the mean participation,
\begin{equation}
  \tilde L_k(f) = L_k(f) + \lambda\, \langle N(f)\rangle_k,
  \label{eq:L_penalized}
\end{equation}
with \(\lambda\) chosen small enough that it only breaks degeneracies between
competing minima.

We consider two competing hypotheses for the target frequency: the full
renormalised Mother Frequency \(F_m^*\) and a halved version \(F_m^*/2\),
motivated by the possibility that the underlying oscillators lock in pairs
rather than individually. In practice this amounts to repeating the
minimisation above with \(F_m^*\) replaced by \(F_m^*/2\). For each family we
retain the combination \((f_{\mathrm{base},k},\text{full or half})\) that
yields the smallest penalised loss \(\tilde L_k\). This procedure is
analogous to model selection between two nested discrete-grainings of the
same underlying frequency scale.

% ----------- IV. RESULTS -----------
\section{Results}
\label{sec:results}

\subsection{Integer participation vs null models}

\figref{fig:N_hist} shows the empirical distribution of
$\Ncorr$ compared with its shuffled-$\ThetaD$ null.
The real data exhibits clear accumulation near small integers
$N \simeq 1$ and $N \simeq 2$ and a pronounced cluster near
$N \simeq 24$--$25$, which is largely absent in the null.
These modes correspond to families whose calibrated $\fbase$
select low-order harmonics.

A more direct comparison is obtained by looking at the residuals
$|\deltalock|$, shown in \figref{fig:delta_hist}.
The real distribution is significantly more concentrated near zero
than the shuffled null:
both KS and MW tests yield $p \lesssim 10^{-13}$,
and Cliff's $\Delta \simeq 0.32$ indicates a medium-to-large effect size.
In other words, after correcting for structural noise,
the dataset contains substantially more near-integer
participation numbers than expected from a smooth frequency spectrum.

\begin{figure}
  \centering
  \includegraphics[width=\linewidth]{figures/fig01a_hist_N_real_vs_shuffle.png}
  \caption{%
  Distribution of noise-corrected participation numbers
  $\Ncorr = \Fmstar / \fbase$ (``Real'') compared with a null model
  obtained by shuffling $\ThetaD$ across materials (``Null'').
  Vertical dotted lines highlight modes around $N \simeq 1$, $N \simeq 2$
  and $N \simeq 24$--$25$.
  }
  \label{fig:N_hist}
\end{figure}

\begin{figure}
  \centering
  \includegraphics[width=\linewidth]{figures/fig01b_hist_delta_real_vs_shuffle.png}
  \caption{%
  Evidence of quantization.%
  \ The distribution of locking residuals
  $|\deltalock| = |\Ncorr - \mathrm{round}(\Ncorr)|$
  for the real data (blue) shows a sharp excess at
  $|\deltalock| \approx 0$ that is absent in the shuffled-$\ThetaD$ null
  (orange).
  Under the null, the probability of obtaining an equal or stronger
  concentration at zero deviation is
  $p < 10^{-13}$ (Kolmogorov--Smirnov and Mann--Whitney tests),
  and Cliff's $\Delta \simeq 0.32$ indicates a medium-to-large effect size.
  }
  \label{fig:delta_hist}
\end{figure}

\subsection{Family-dependent locking strength}

The locking residuals vary systematically across families.
\figref{fig:delta_by_family} summarizes the distribution of
$|\deltalock|$ (clipped at 1.0 for readability) for each category.
Iron-based superconductors and heavy-fermion systems show
the smallest median residuals, indicating strong integer participation,
whereas type-I and high-pressure materials exhibit broader tails.

\begin{figure*}
  \centering
  \includegraphics[width=\linewidth]{figures/fig02_delta_by_family.png}
  \caption{%
  Per-family distribution of locking residuals $|\deltalock|$
  (clipped at 1.0).
  Boxes show interquartile ranges and medians; whiskers extend
  to $1.5$ times the interquartile range.
  Families are ordered by increasing median residual.
  }
  \label{fig:delta_by_family}
\end{figure*}

\subsection{Structural noise vs integer locking}

We now relate locking strength to structural noise.
\figref{fig:delta_vs_noise} plots $|\deltalock^{(i)}|$
against the noise $z$--score $\znoise^{(i)}$, with points
colored by family.
Despite considerable scatter, there is a clear positive
monotonic trend:
the Spearman correlation is
$\rho \simeq 0.34$ with $p \simeq 5.6 \times 10^{-4}$.
Materials with higher structural noise tend to deviate more from
integer participation.

\begin{figure}
  \centering
  \includegraphics[width=\linewidth]{figures/fig03a_delta_vs_noise_scatter.png}
  \caption{%
  Absolute locking residual $|\deltalock|$ vs structural-noise
  $z$--score $\znoise$.
  Colors indicate material families.
  The Spearman correlation is $\rho \simeq 0.34$
  with $p \simeq 5.6 \times 10^{-4}$,
  indicating that higher structural noise is associated with weaker
  integer locking.
  }
  \label{fig:delta_vs_noise}
\end{figure}

The same relationship is visible when comparing groups.
\figref{fig:noise_almost_integer} shows the distribution of
$\znoise$ for the ``almost-integer'' group
(top 20\% lowest $|\deltalock|$) and the rest of the dataset.
The almost-integer materials have markedly lower noise,
with Cliff's $\Delta \simeq -0.47$, a large negative effect size.

\begin{figure}
  \centering
  \includegraphics[width=\linewidth]{figures/fig03b_noise_almost_integer_vs_rest.png}
  \caption{%
  Structural-noise $z$--score $\znoise$ for materials in the
  ``almost-integer'' group (top 20\% smallest $|\deltalock|$)
  and for the rest of the dataset.
  The almost-integer group is shifted toward lower noise,
  with Cliff's $\Delta \simeq -0.47$.
  }
  \label{fig:noise_almost_integer}
\end{figure}

\subsection{Base-frequency hierarchy across families}

The calibrated base frequencies $\fbase^{(k)}$
form a simple, interpretable pattern.
\figref{fig:fbase_by_family} shows the median $\fbase^{(k)}$
per family, ordered by increasing value, together with the
global median $\fbase_{\mathrm{med}} \simeq 4.157$.

Iron-based superconductors cluster near
$\fbase \simeq 1$, suggesting that, in these systems, the mother
frequency often participates in the fundamental mode and that
$\Fmstar$ itself already behaves as an integer participation number.
Heavy-fermion and oxide materials sit near
$\fbase \simeq 2$, consistent with a first harmonic.
Binary, molecular and high-pressure superconductors, as well as
type-I materials, concentrate around $\fbase \simeq 4$,
the same integer that controls the superfluid--roton ratio in \HeFour.
The independently calibrated $\fbase$ for the Superfluid family
is $\simeq 4.0$, so the global median effectively recovers the
\HeFour\ scale when pooling all families.

\begin{figure*}
  \centering
  \includegraphics[width=\linewidth]{figures/fig04a_fbase_by_family.png}
  \caption{%
  Base frequency $\fbase$ per family (median over materials), ordered by
  increasing value.
  The dashed line marks the global median
  $\fbase_{\mathrm{med}} \simeq 4.157$.
  Families cluster around approximately harmonic values
  $\fbase \sim 1$, $2$ and $4$, with iron-based superconductors
  near the fundamental and binary / molecular / high-pressure
  materials near the $n=4$ harmonic associated with \HeFour.
  }
  \label{fig:fbase_by_family}
\end{figure*}

% ----------- V. DISCUSSION -----------
\section{Discussion and outlook}\label{sec:discussion}

The first study in this programme showed that a delayed-oscillator locking
grammar on the primes \(\{2,3,5,7\}\), combined with a universal correction
law calibrated on classical metals, is sufficient to compress a heterogeneous
set of superconductors and superfluid helium into stable prime-space
fingerprints. The present work adds a second layer: an integer–participation
structure in the renormalised Mother Frequency, once structural noise is
taken into account.

On the phenomenological side, the main result is that for each macroscopic
family there exists a simple base frequency \(f_{\mathrm{base},k}\) such that
the renormalised ratios \(F_m^*/f_{\mathrm{base},k}\) cluster tightly around
integers. This clustering is robust under cross-validation and against
multiple null models, and it correlates systematically with a structural
noise field \(\hat\xi\) derived from independent considerations of cluster
topology and vector dissipation. Materials that sit closest to an integer
participation number tend to exhibit lower structural noise, while strongly
non-integer cases concentrate in families where geometric frustration or
complex topology are expected.

From a DOFT-inspired viewpoint, these findings are consistent with a picture
in which superconducting clusters behave as discrete assemblies of delayed
oscillators that seek integer harmonics of a family-specific base frequency,
with structural noise quantifying the misfit between the actual geometry and
an idealised closed configuration. From a more conventional perspective, the
integer structure may simply reflect a coarse-grained count of effective
degrees of freedom (layers, chains, or coupled sub-networks) that contribute
coherently to the superconducting state.

We deliberately remain agnostic about the microscopic origin of
\(f_{\mathrm{base},k}\). One speculative route is to connect it to known
characteristic frequencies in each family, such as dominant phonon branches,
spin fluctuations, or collective modes observed in Raman and neutron
scattering. Another is to interpret the integer participation as a proxy for
cluster geometry, in analogy with magic numbers in metallic clusters, where
closed icosahedral shells lead to enhanced stability. At present, the data
are not sufficient to distinguish between these scenarios.

The comparison with superfluid \(^4\)He should also be treated with care.
The appearance of a low integer scale linked to the roton minimum is
intriguing, but bosonic superfluidity and fermionic superconductivity rely
on different microscopic mechanisms. For now we view the helium case as an
interesting boundary point for the locking grammar rather than as evidence
for a unified oscillator mechanism.

Several extensions are natural. On the empirical side, a larger and more
balanced dataset per family would allow sharper estimates of
\(f_{\mathrm{base},k}\), more stringent cross-validation, and a clearer
separation between bulk-like and strongly frustrated regimes. On the
theoretical side, it would be valuable to embed the integer participation
directly into a microscopic delayed-oscillator field theory, or to derive
effective \(f_{\mathrm{base},k}\) from simplified models of coupled phonons
and electrons. Ultimately, the usefulness of the present construction will
be decided by its predictive power: whether integer participation, combined
with the prime-space fingerprints, can help identify promising new
superconductors or clarify why certain families systematically fall short of
their apparent locking potential.
%\section{Discussion}
%\label{sec:discussion}
%
%The analysis above establishes three empirical facts:
%
%(i) After correcting for structural noise, the noise-corrected
%frequencies $\Fmstar$ admit a family-dependent base frequency
%$\fbase^{(k)}$ such that the rescaled quantities $\Ncorr$
%cluster unusually close to integers.
%
%(ii) The strength of this clustering, quantified by $|\deltalock|$,
%is modulated by structural noise both continuously
%(\figref{fig:delta_vs_noise}) and in group comparisons
%(\figref{fig:noise_almost_integer}),
%with low-noise materials displaying the most pronounced integer locking.
%
%(iii) The calibrated base frequencies fall into a simple hierarchy
%of approximate harmonics: $\fbase \sim 1$ for iron-based,
%$\sim 2$ for oxides and heavy fermions, and $\sim 4$ for binary,
%molecular, high-pressure and type-I superconductors, with a global
%median $\simeq 4.16$ very close to the $n=4$ ratio
%characteristic of the roton branch in \HeFour.
%
%One interpretation is that the mother frequency $\Fm$
%is not an arbitrary thermal ratio but a coarse probe of an underlying
%oscillator network whose normal modes are discretized.
%Structural noise---arising from defects, frictional heating and
%pressure-induced distortions---broadens and shifts these modes.
%Where noise is low, the projection of $\Fmstar$ onto this discrete
%spectrum becomes visible as integer participation.
%Where noise is high, the same structure is smeared into the null.
%
%We emphasize that our analysis is agnostic about the microscopic
%origin of these oscillators.
%They could arise from phase-coherent phonon modes, charge or spin
%density waves, or more exotic collective excitations.
%Our results do not prove a specific mechanism;
%they only constrain any viable mechanism to reproduce the observed,
%family-dependent integer structure once mapped into $\Fm$.
%
%At the same time, the numerical coincidence between the global
%$\fbase$ median and the \HeFour\ roton ratio suggests that
%superconductivity and superfluidity may share a common ``unit''
%of coherence, manifested at the level of a simple frequency ratio.
%Whether this reflects a deep universality class or a more mundane
%convergence of scales remains an open question.

% ----------- VI. CONCLUSION -----------
\section{Conclusion}
\label{sec:conclusion}

We have shown that a simple frequency ratio
$\Fm = \ThetaD / \Tc$, when corrected for structural noise
and rescaled by a family-dependent base frequency,
exhibits robust integer participation across a heterogeneous set
of superconductors and liquid \HeFour.
The effect is statistically strong, structurally modulated and
organized by approximate harmonics of a base frequency whose global
scale matches that of \HeFour.

Among all families, iron-based superconductors stand out as the
clearest realization of this picture: with $\fbase \simeq 1$,
they operate in a natural resonance unit where the noise-corrected
mother frequency $\Fmstar$ itself already acts as an integer
participation number.
In this sense, they behave as ideal, unscaled harmonic oscillators
within the phenomenological framework developed here.

Rather than offering a full microscopic theory, our goal is to
provide a compact phenomenological summary of patterns that any
future theory ought to explain.
The framework---mother frequency, structural noise, base-frequency
calibration and integer participation---is deliberately minimal
and can be applied to expanded datasets as they become available.

%------------------------------------------------------------
\section*{Acknowledgements}

The author thanks colleagues and online communities for discussions and for
making high-quality superconductivity and superfluid datasets publicly
available. This project made extensive use of open-source software, including
Python, NumPy, SciPy, pandas, Matplotlib, and Jupyter.

AI tools were used as assistants in the development of the code and this
manuscript. In particular, OpenAI GPT-5.1~Thinking, Google Gemini~Pro, and
Anthropic Claude were employed for drafting text, refactoring analysis
pipelines, and checking statistical procedures. All modelling decisions,
dataset curation and interpretations presented here were made and verified by
the human author.

% ----------- APPENDIX: SUPPLEMENTARY FIGURES -----------
\clearpage
\appendix

\onecolumngrid 

\section*{Appendix: Structural-noise diagnostics}

\begin{figure}[h!]
  \centering
  \begin{minipage}{0.48\textwidth}
    \centering
    \includegraphics[width=\linewidth]{figures/figS01_noise_by_family.png}
    \caption{%
    Structural-noise $z$--score $\znoise$ by family
    (only families with a non-degenerate noise model are shown).
    Distributions are standardized within each family;
    boxes show interquartile ranges and medians.
    }
    \label{fig:noise_by_family}
  \end{minipage}\hfill 
  \begin{minipage}{0.48\textwidth}
    \centering
    \includegraphics[width=\linewidth]{figures/figS02_Tc_vs_Tc_ideal.png}
    \caption{%
    Measured transition temperature $\Tc$ vs ideal temperature
    $\Tcone = \Tc (1 + \widehat{\xi})$.
    The dashed line indicates $\Tcone = \Tc$.
    Most materials lie close to the identity, with a subset
    receiving larger upward corrections due to structural noise.
    }
    \label{fig:Tc_vs_Tc_ideal}
  \end{minipage}
\end{figure}
\twocolumngrid

% ----------- Bibliography -----------
\nocite{*}
\bibliographystyle{apsrev4-2}
\bibliography{main}

\end{document}
